% Solutions/hints for ch05 exercises (assembled into Appendix C)

\begin{solution}{exr:ch05-states}
(a) $\rhs(s)$ rises to the new value (unless a second predecessor offers the old total, in which case nothing changes) while $\gcost(s)$ keeps the old, smaller value: $s$ is underconsistent with key $[\gcost(s)+\hcost(s);\,\gcost(s)]$, built from the old value. (b) $\rhs(s)$ drops below $\gcost(s)$: overconsistent, key $[\rhs(s)+\hcost(s);\,\rhs(s)]$, built from the new value. (c) $\rhs(s)$ is a minimum over the edges \emph{into} $s$, so it does not change and $s$ stays consistent; it is the head of the changed edge that \LpaUpdateVertex is called on and that may become underconsistent. (d) The minimum in \cref{eq:ch05-rhs} is attained by the second predecessor with the same total, so $\rhs(s)$ and hence the state of $s$ do not change; \LpaUpdateVertex still runs on $s$ (it is a successor of the popped vertex) but removes and re-inserts nothing.
\end{solution}

\begin{solution}{exr:ch05-handtrace}
Blocking $(1,0)$ changes the eight edges around it, so \LpaUpdateVertex runs on $(1,0)$, $(0,0)$, $(2,0)$ and $(1,1)$. Only $(1,0)$ changes: its $\rhs$ was $2$ through $(1,1)$ and becomes $\infty$, which equals its $\gcost$, so it simply leaves the queue. $(0,0)$ keeps $\rhs=1$ through $S$, $(2,0)$ keeps $\rhs=3$ through $(2,1)$, and $(1,1)$ keeps $\rhs=1$; none of them ever relied on $(1,0)$, whose $\gcost$ was still $\infty$. The queue now holds nine cells with keys $[6;1],[6;1],[6;2],[6;3],[6;3],[6;4],[6;4],[6;5],[6;5]$, the goal is consistent with $\gcost(T)=4$, and the smallest key $[6;1]$ is not smaller than $k(T)=[4;4]$: \LpaComputeShortestPath returns after zero expansions and the path along the middle row stands. The difference to \cref{tab:ch05-lpastar-repair} is that $(3,1)$ carried a finite $\gcost$-value on which the goal depended, whereas $(1,0)$ was never expanded, so no stored value was stale. Unblocking $(3,1)$ in the state of the right panel calls \LpaUpdateVertex on $(3,1)$ and its four neighbours; $\rhs(3,1)$ falls from $\infty$ to $3$ through $(2,1)$, so $(3,1)$ is \emph{over}consistent with key $[4;3]$, while the neighbours are unaffected because $\gcost(3,1)$ is still $\infty$. The repair pops $(3,1)$ (key $[4;3]$, $\gcost\gets 3$), which lowers $\rhs(T)$ from $6$ to $4$ and queues $T$ with key $[4;4]$, then pops $T$ ($\gcost\gets 4$); two expansions, and the path is the middle row again. A cost decrease can only lower $\rhs$-values, so only overconsistent vertices appear, and the queued side cells with keys $[6;\cdot]$ are never touched.
\end{solution}

\begin{solution}{exr:ch05-secondkey}
Take the vertices $S,u,w,T$ with edges $S\to u$ (cost $1$), $u\to w$ ($1$), $w\to T$ ($1$) and $u\to T$ ($2$), start $S$, goal $T$, and the exact heuristic $\hcost(S)=3$, $\hcost(u)=2$, $\hcost(w)=1$, $\hcost(T)=0$, which stays consistent when costs increase. Both routes to $T$ cost $3$, so every vertex has $k_1=3$ in the first search. With $k_1$ only and ties broken against us, $T$ is expanded before $w$, and $w$ stays in the queue with $\gcost(w)=\infty$ and $\rhs(w)=2$ derived from $\gcost(u)=1$. Now raise $c(S,u)$ to $2$ and $c(u,T)$ to $10$ in one batch. \LpaUpdateVertex gives $\rhs(u)=2$ ($u$ underconsistent, $k_1=1+2=3$) and $\rhs(T)=\min(1+10,\infty)=11$ ($T$ underconsistent, $k_1=3$); $w$ keeps $k_1=2+1=3$. Pop $T$ (underconsistent, $\gcost(T)\gets\infty$, re-queued with $\rhs=11$), then $w$ \emph{first} among the ties: it is overconsistent and accepts the stale value $\gcost(w)=2$, which makes $\rhs(T)=3$ and $k_1(T)=3$. Pop $u$: $\gcost(u)\gets\infty$, so $\rhs(w)=\infty$ and $w$ is underconsistent with $k_1=3$, while $u$ re-enters with $\rhs(u)=2$, $k_1=4$. Pop $w$ a second time ($\gcost(w)\gets\infty$; $T$ and $w$ both become consistent at $\infty$ and leave the queue), pop $u$ ($\gcost(u)\gets 2$, $\rhs(w)=3$, $\rhs(T)=12$), pop $w$ a \emph{third} time ($\gcost(w)\gets 3$, $\rhs(T)=4$) and finally $T$ with $\gcost(T)=4$: seven expansions and $w$ three times. If instead the tie after the second step is broken toward $T$, the loop even stops with the wrong value $\gcost(T)=3$. With the two-component key the first search expands $w$ ($[3;2]$) before $T$ ($[3;3]$), and after the change the pops are $u\,[3;1]$, $w\,[3;2]$, $T\,[3;3]$ (all underconsistent), then $u\,[4;2]$, $w\,[4;3]$, $T\,[4;4]$: six expansions, each vertex twice, and the correct cost $4$.
\end{solution}

\begin{solution}{exr:ch05-kmproof}
Let $s$ be in $U$ with the key stored when the robot stood at $s_{\mathrm{last}}$ and the modifier had the value $k_m$: $k_1^{\mathrm{old}}=m+\hcost(s_{\mathrm{last}},s)+k_m$ with $m=\min(\gcost(s),\rhs(s))$, and $k_2^{\mathrm{old}}=m$. Neither $\gcost(s)$ nor $\rhs(s)$ changed since, because any change goes through \DslUpdateVertex, which re-inserts $s$ with a fresh key. After line~\ref{alg:ch05-dstarlite:km} the current key is $k_1=m+\hcost(s_{\mathrm{start}},s)+k_m+\hcost(s_{\mathrm{last}},s_{\mathrm{start}})$, and the triangle inequality $\hcost(s_{\mathrm{last}},s)\le\hcost(s_{\mathrm{last}},s_{\mathrm{start}})+\hcost(s_{\mathrm{start}},s)$ gives $k_1\ge k_1^{\mathrm{old}}$ with $k_2=k_2^{\mathrm{old}}$, so the stored key is a lower bound. If the batch is the second one since the key was stored, apply the argument twice: the sum of the two legs is at least the direct distance, again by the triangle inequality. For the counterexample take \cref{tab:ch05-km}: $u$ was stored with $[10;4]$; the robot moves two cells toward $u$ and, without the modifier, the current key is $[8;4]$, so the stored key is a strict upper bound. A vertex $v$ inserted now with $[9;6]$ is popped before $u$ although its current key is larger, which breaks invariant~(iii) of the proof sketch of \cref{thm:ch05-correct}, that vertices are expanded in nondecreasing order of their current keys; $v$ may then be expanded with an $\rhs$-value that still depends on $u$, and the loop may stop with $u$ buried under an inflated key while $\gcost(s_{\mathrm{start}})$ is not the true distance. Finally, $k_m$ must be accumulated: keys stored at different times were computed with different values of $k_m$ and different robot positions, and each of them needs the offset of the legs walked \emph{since its insertion}. The sum of the legs is at least $\hcost(s_{\mathrm{last}},s_{\mathrm{start}})$ for every one of them, whereas the direct distance $\hcost(s_{\mathrm{orig}},s_{\mathrm{start}})$ from the original start can shrink when the robot turns back (it is $0$ if the robot returns to $s_{\mathrm{orig}}$), and then keys computed later would be smaller than keys stored earlier, which is the upper-bound situation above.
\end{solution}

\begin{solution}{exr:ch05-directions}
Hint and sketch. In the reversed graph $G^R$ one has $\Pred_R(s)=\Succ(s)$ and $c^R(s',s)=c(s,s')$, so \lpastar's lookahead \cref{eq:ch05-rhs} on $G^R$ with the roles of start and goal swapped becomes $\rhs(s)=\min_{s'\in\Succ(s)}\bigl(c(s,s')+\gcost(s')\bigr)$, which is line~\ref{alg:ch05-dstarlite:rhs} of \cref{alg:ch05-dstarlite}, and the seed $\rhs(s_{\mathrm{goal}})=0$ is line~\ref{alg:ch05-dstarlite:init}. The query vertex of the reversed search is $s_{\mathrm{start}}$, so the heuristic must be measured from it, $\hcost(s_{\mathrm{start}},\cdot)$, and the loop condition becomes $\IncTopKey{} < k(s_{\mathrm{start}})$ or $\rhs(s_{\mathrm{start}})\neq\gcost(s_{\mathrm{start}})$ (line~\ref{alg:ch05-dstarlite:while}); an expansion updates $\Succ_R(u)=\Pred(u)$ (lines~\ref{alg:ch05-dstarlite:over}--\ref{alg:ch05-dstarlite:underpred}), and the path is read forward from $s_{\mathrm{start}}$ over the successor minimising $c+\gcost$. That is exactly \cref{alg:ch05-dstarlite} with $k_m\equiv 0$ and with lines~\ref{alg:ch05-dstarlite:kold}--\ref{alg:ch05-dstarlite:reinsert} deleted. Only four things are there because the start moves: line~\ref{alg:ch05-dstarlite:move} (the robot step, which reassigns $s_{\mathrm{start}}$), line~\ref{alg:ch05-dstarlite:km} together with $s_{\mathrm{last}}\gets s_{\mathrm{start}}$ (the key modifier), and the $k_{\mathrm{old}}$ test of lines~\ref{alg:ch05-dstarlite:kold}--\ref{alg:ch05-dstarlite:reinsert}, which repairs stored keys that the moves have turned into strict lower bounds. Everything else is \lpastar read backwards.
\end{solution}

\begin{solution}{exr:ch05-coding}
Expected outcome. (a) The two path costs must agree at \emph{every} tick of every run;
a single disagreement means a bug, and the usual causes are calling \DslUpdateVertex on
the head instead of the tail of a changed edge, forgetting $\Pred(u)\cup\set{u}$ in the
underconsistent case, or updating $k_m$ once per changed cell instead of once per batch.
(b) Most events change a cell that is not on the current path and not in the expanded
region, and \dstarlite should then cost $0$ expansions, while \astar re-expands its
whole search region; over a run of a few dozen events expect \dstarlite to expand an
order of magnitude fewer vertices in total, as in \cref{fig:ch05-experiment}. The first
search is the exception: it is several times more expensive than \astar (a factor of
about $4.6$ in \cref{sec:ch05-experiment}), because the second key component expands the
whole $\fcost$-band, so plot event $0$ separately or the axis will hide everything else.
Wall-clock time need not follow the expansion count: the per-expansion cost of
\dstarlite is several times that of \astar, so on small grids repeated \astar can stay
ahead in time while losing badly in expansions.
\end{solution}

\begin{solution}{exr:ch05-threshold}
Expected outcome. For very small $\phi$ (a handful of changed cells) \dstarlite wins by
orders of magnitude in \emph{expansions}: most changes touch cells whose goal distance
does not change at all, and the repair costs almost nothing, while \astar re-expands its
whole search region every time. The crossover in \emph{time} comes much earlier than the
crossover in expansions, because one \dstarlite expansion costs several times an \astar
expansion --- about $28\,\mu$s against $4\,\mu$s in \cref{sec:ch05-experiment} ---
so \dstarlite can already be losing in seconds while still winning by a factor of ten in
expansions. Around a $\phi$ of a few per cent on a $100\times 100$ grid the changes are
no longer local: the repair touches a constant fraction of the vertices,
\cref{thm:ch05-expansions} then permits up to $2|V|$ expansions plus a
\DslUpdateVertex on every neighbour of each of them, and a fresh \astar --- which
expands every vertex at most once and relaxes each edge in $\bigO{1}$ --- is faster. At
$\phi=0.5$ nothing of the old search survives and \dstarlite is strictly worse.
The rule for the replanning layer of \cref{ch:ch24} is therefore: count the changed
cells reported per replanning cycle, keep a running estimate of the two costs, and fall
back to \astar from scratch above the measured threshold --- and always after a full map
replacement, a goal change, or a re-localisation that shifts the whole grid, since none
of the stored values is then worth repairing. Exact numbers are implementation- and
machine-dependent (a compiled indexed heap moves the crossover to the right, an
interpreted lazy queue to the left); report your own crossover, in expansions and in
time, and say which grid and which language produced it.
\end{solution}

\begin{solution}{exr:ch05-intruder}
Hint and sketch. (a) For each tick $t\le H$ rasterise the disc of radius $r_t=r_0+\sigma t$ around $\pos_t$, inflated by the drone's own radius, and take $B=\bigcup_t B_t$; raise the cost of every edge into a cell of $B$ to $\infty$ and call \DslUpdateVertex on each such cell and its neighbours, as in \cref{sec:ch05-implementation}. When the next prediction arrives, do \emph{not} rebuild $B$: form the new set $B'$, and call \DslUpdateVertex only on the symmetric difference $B\triangle B'$ (newly blocked cells, and cells released because the intruder has passed or the prediction shifted) and their neighbours. Cells in $B\cap B'$ cost nothing, which is the whole point of an incremental search; keep $B$ between calls so the difference can be formed in $O(|B|+|B'|)$. (b) Blocking the whole tube for the whole horizon is conservative in space \emph{and} in time: a cell that the intruder occupies for two ticks is forbidden for all $H$, so a corridor that would be free by the time the drone reaches it is closed, and the planner either takes a long detour or reports \emph{no path} although a safe schedule exists. A space-time search (\cref{ch:ch04}) over states $(\text{cell},t)$ blocks $(\text{cell},t)$ only for the ticks in which $B_t$ covers the cell; edge changes then touch a single time layer, so the repair is even more local, at the price of a state space $H$ times larger and of a plan that must be re-timed rather than just re-routed. (c) If the goal moves, the $\gcost$-values, which are distances \emph{to the goal}, change with every goal step, and the property that made the backward search cheap is gone; $k_m$ cannot help, because the goal enters the $\gcost$-values and not the heuristic. One must either re-initialise or use Moving Target \dstarlite \cite{sun2010moving}, which salvages the still-valid part of the search tree and discards the rest.
\end{solution}
